The \aschip{} implements a minimal \textbf{single external interrupt source,
machine-mode-only} interrupt concept, identical in both chip variants
(\texttt{as\_cpux} in \texttt{RV\_NoPipeline} and \texttt{cc\_cpupipefwd} in
\texttt{RV\_PipelineCacheFwd} implement the same CSR/trap logic).
There is no timer interrupt (no \texttt{mtime}/\texttt{mtimecmp}) and no
software interrupt (\texttt{MSIP}); only the RISC-V \textbf{Machine External
Interrupt} (MEI, \texttt{mcause}\,=\,\texttt{0x800...B}, cause code 11) is
implemented.

\paragraph{Peripheral-to-CPU path.}
The top level (\texttt{as\_top\_mem}) presents an 8-bit vector,
\texttt{irq\_ext\_i} (\texttt{irq\_total\_num\_ext\_c}\,=\,8), to the CPU.
Of these 8 bits, \textbf{only bit~7 is actually wired into the CPU's
interrupt logic}: the core synchronises \texttt{irq\_ext\_i[7]} through a
2-flip-flop synchroniser and feeds it directly into \texttt{mip.MEIP}.
Bits~\texttt{[6:0]} are declared but not consumed by either CPU
implementation (marked \texttt{UNUSEDSIGNAL} in \texttt{asCPUx.sv}).
Bit~7 is driven by \textbf{ICU0} (Section~\ref{sec:icufunc01}), a
Wishbone-addressable interrupt controller that aggregates GPIO0, UART0,
and (on \texttt{RV\_PipelineCacheFwd}) QSPI0 with software-configurable
priority and per-source masking, plus 5 further source slots reserved for
future peripherals (Table~\ref{tab:asIcuSrcAssign}). GPIO0's own interrupt
chain (per-pin change detection $\rightarrow$ \texttt{IRQSS}/\texttt{IRQSM}
$\rightarrow$ single-bit \texttt{RIS[0]} $\rightarrow$ \texttt{IMSC[0]}
$\rightarrow$ \texttt{MIS[0]} $=$ \texttt{gpio\_irq\_o}) collapses all GPIO
pins onto one line before it ever reaches ICU0; see
Section~\ref{sec:gpiofunc01} for the full GPIO0 register description.
On \texttt{RV\_NoPipeline}, UART0's interrupt now reaches the CPU for the
first time via ICU0 -- previously it was wired directly to
\texttt{irq\_ext\_i[6]}, a bit the CPU never reads, so it had no effect on
program flow.

\paragraph{CPU-side CSR mechanism.}
\begin{itemize}
  \item \texttt{mip[11]} (MEIP) is \textbf{read-only}: it directly mirrors
    the synchronised \texttt{icu\_irq\_o} line and cannot be set or
    cleared by software: it stays 1 for as long as ICU0 has at least one
    pending, unmasked, globally-enabled source, and clears only when
    firmware clears the responsible source(s) via ICU0's own
    \textbf{ICU0\_ICR} (and, typically, the underlying peripheral's own
    clear register too).
  \item \texttt{mie[11]} (MEIE) is a normal read/write CSR bit, enabled
    (\texttt{=1}) by default out of reset.
  \item An interrupt is taken when
    \texttt{mip[11]\,\&\,mie[11]\,\&\,mstatus.MIE} all hold.
  \item \texttt{mtvec} supports \textbf{Direct mode only}
    (\texttt{mtvec[1:0]}\,=\,\texttt{00}); every trap -- exception or
    interrupt alike -- jumps to the same handler address, reset default
    \texttt{0x0000\_0000\_0000\_7F00}.
    There is no vectored mode and no per-cause offset.
  \item On trap entry: \texttt{mepc} $\leftarrow$ PC (faulting instruction,
    for illegal-instruction/misaligned-address/\texttt{ecall}/\texttt{ebreak}
    exceptions) or PC$+4$ (return address, for the external interrupt);
    \texttt{mstatus.MPIE}
    $\leftarrow$ \texttt{mstatus.MIE}; \texttt{mstatus.MIE} $\leftarrow$ 0;
    PC $\leftarrow$ \texttt{mtvec}.
  \item \texttt{mret} restores \texttt{mstatus.MIE} $\leftarrow$
    \texttt{mstatus.MPIE} and sets \texttt{mstatus.MPIE} $\leftarrow$ 1.
  \item \textbf{Trap priority} (highest first): illegal instruction
    $>$ misaligned load/store address $>$ \texttt{ecall}/\texttt{ebreak}
    $>$ external interrupt.
    Only one trap is taken per cycle; a pending external interrupt is
    deferred behind any pending exception.
    In practice this ordering is never contended: \texttt{ecall}/\texttt{ebreak}
    are a single, unambiguous instruction encoding that cannot simultaneously
    be illegal or trigger a misaligned load/store.
\end{itemize}

\paragraph{Synchronous exceptions: \texttt{ecall} and \texttt{ebreak}.}
Besides the illegal-instruction and misaligned-address traps, both CPU
variants decode the two mandatory RV64I \texttt{SYSTEM} exceptions directly:
\begin{itemize}
  \item \texttt{ecall} (\texttt{funct3}\,=\,\texttt{000},
    \texttt{imm[11:0]}\,=\,\texttt{0x000}) sets
    \texttt{mcause}\,=\,11 (\textit{Environment call from M-mode}).
  \item \texttt{ebreak} (\texttt{funct3}\,=\,\texttt{000},
    \texttt{imm[11:0]}\,=\,\texttt{0x001}) sets
    \texttt{mcause}\,=\,3 (\textit{Breakpoint}).
\end{itemize}
Both traps behave like the illegal-instruction/misaligned-address case:
\texttt{mepc} points at the \texttt{ecall}/\texttt{ebreak} instruction itself
(not PC$+4$), \texttt{mstatus.MPIE}/\texttt{MIE} update the same way, and
\texttt{mret} returns control to the trap handler's chosen resume address.
Neither instruction has any side effect beyond the trap -- there is no
environment-call ABI and no hardware breakpoint comparator; both exist purely
to enter the trap handler, which is what \texttt{riscv-arch-test} compliance
and a future GDB-stub breakpoint mechanism require. In
\texttt{RV\_NoPipeline} (\texttt{as\_cpux}), the redirect to \texttt{mtvec}
is triggered one cycle early via \texttt{ecbrk\_pending\_s} against the
live I-Mem prefetch, mirroring how \texttt{mret\_pending\_s} already handles
\texttt{mret} -- otherwise the design's always-ack instruction prefetch would
transiently execute one stale, wrong-path instruction before the redirect
lands. \texttt{RV\_PipelineCacheFwd} (\texttt{cc\_cpupipefwd}) needs no such
trick: its I-Cache uses an explicit req/rvalid handshake, so the ordinary
commit-stage flush already used for illegal-instruction/misaligned traps is
sufficient.

\paragraph{Open items.}
ICU0 provisions 8 source slots but only 3 are wired to real hardware on
either chip (GPIO0, UART0, and -- \texttt{RV\_PipelineCacheFwd} only --
QSPI0); connecting a future interrupt-capable peripheral only requires
wiring it into one of the 5 remaining slots
(Table~\ref{tab:asIcuSrcAssign}), not a redesign of the register map.
Both CPUs still support \texttt{mtvec} Direct mode only -- ICU0's
\textbf{ICU0\_IRQID} register lets firmware branch on cause in software
from the one shared handler, but hardware-vectored dispatch (a per-cause
jump table) would require CPU changes outside ICU0's scope and remains
unimplemented. Nested/preemptive interrupt handling (a higher-priority
interrupt preempting an in-progress lower-priority ISR) is likewise not
implemented; RISC-V M-mode already precludes it at the CPU level
(\texttt{mstatus.MIE} is cleared on trap entry) independent of ICU0.